This class encapsulates the information required to determine an
effective address at runtime. The general
representation for an address is a sum of two registers and a
constant; this may change in future releases. Some
architectures use only certain bits of a register (e.g. bits 25:31 of
XER register on the Power chip family); these are
represented as pseudo-registers. The numbering scheme for registers
and pseudo-registers is implementation dependent
and should not be relied upon; it may change in future releases.

\begin{apient}
  int getImm()
\end{apient}

\apidesc{
  Return the constant offset. This may be positive or negative.
}

\begin{apient}
  int getReg(unsigned i)
\end{apient}

\apidesc{
  Return the register number for the ith register in the sum, where 0
  ${\leq}$ i ${\leq}$ 2. Register numbers are
  positive; a value of -1 means no register.
}

\begin{apient}
  int getScale()
\end{apient}

\apidesc{
  Returns any scaling factor used in the memory address computation.
}
